\honest{Reductionism}{Eliezer Yudkowsky}{2008}

In April 2007 Matthew C asked me to explain ``How and why the current reigning philosophical hegemon (reductionistic materialism) is obviously correct'' while the views of all past civilizations are suspect. I deemed the request legitimate, but could not take it up before the Mind Projection Fallacy posts. \nb{The full comment reads as sarcasm. It ends by proposing ``some well-placed rocks thrown at anyone who would challenge the current paradigm, since they are obviously just fools or liars.''}

Leaving materialism for later, I hold that reductionism, in my sense, ``is obviously correct; and to perdition with any past civilizations that disagreed.'' That sounds strong. Even General Relativity may be overturned. But science does not go back to a theory as refuted as Newtonian mechanics. The ratchet does not turn in reverse, so dismissing past civilizations is safe when what they believed has been falsified. \nb{The post does not say which past view about reduction has been falsified.} Besides, reductionism is not so much a positive hypothesis as ``the absence of belief''.

A man who said he had been a Navy gunner told me that computing a shell's trajectory with relativity gives the wrong answer. I said no: relativity might be too slow to compute in time, but its answer is always more accurate. He insisted that things at a shell's speed are governed by Newtonian mechanics, not relativity. If that were true, I said, he could collect a Nobel Prize. \nb{At a shell's speed, relativity's corrections are of order a billionth or less; small, but they are corrections.}

Physics uses one fundamental theory for a 747 and for gold nuclei colliding: ``special relativity, quantum mechanics, and chromodynamics.'' \nb{Chromodynamics is the theory of the strong force. What governs the airplane's aerodynamics is electromagnetism, and the fundamental theory is the Standard Model. The argument does not depend on the slip.} But we model the 747 with wings and airflow, because a fundamental model would take ``a gazillion years'' to run and would not fit on all the world's computers. A program modeling the 747's aerodynamics may not have a single bit that stands for a quark. Is the 747 made of something other than quarks? No. The model's elements simply do not correspond one to one with the quarks. You cannot fold the territory into your glove compartment, and a map's scale is a fact about the map. A fundamental model, if it could run, would predict better than the aerodynamic one, and need contain nothing explicit about wings or lift. \nb{Anderson's ``More Is Different'' (1972) made this distinction: reducing everything to fundamental laws ``does not imply the ability to start from those laws and reconstruct the universe.''}

``What?'' cries the antireductionist. ``Are you telling me the 747 doesn't really have wings?'' \nb{This imagined voice and the gunner are the only opponents the post answers.} My answer is subtler than saying an object has different descriptions at different levels: having different descriptions at different levels belongs to talk about maps, not talk about territory. The airplane and the laws of physics do not use levels, as the gunner thought. We use them, for convenience. A fundamental model would contain all the facts about the wings, implicitly. You, studying it, could work out where the wings are, and then the wings would be explicit, in your mind.

The levels feel real because that is how an efficient multi-level model feels from inside, just as a belief feels like looking straight at reality. The brain compresses an object it cannot begin to model fundamentally. The airplane, and even a hydrogen atom, is too large to model that way. You can't handle the truth, but reality handles it with no simplification (I wish I knew where it gets the computing power).

``The way physics really works, as far as we can tell, is that there is only the most basic level.'' The laws contain no extra causal entities for lift or wings, as an engineer's mind contains extra mental entities for them. Reductionism ``is not a positive belief, but rather, a disbelief that the higher levels of simplified multilevel models are out there in the territory.'' \nb{The two sentences state one thesis, positively and negatively.} This dissolves the question of the wings. ``The critical words are really and see.''
